% BHU PYQ -- Professional Exam Template
\documentclass[11pt,a4paper]{article}

\usepackage[a4paper,top=2.5cm,bottom=2.5cm,left=2.8cm,right=2.8cm,headheight=14pt]{geometry}
\usepackage{amsmath,amssymb}
\usepackage[T1]{fontenc}
\usepackage[utf8]{inputenc}
\usepackage{lmodern}
\usepackage{microtype}
\usepackage{fancyhdr}
\usepackage{enumitem}
\usepackage{listings}
\usepackage{hyperref}
\usepackage{lastpage}
\usepackage{parskip}

\hypersetup{colorlinks=true,urlcolor=black,linkcolor=black,
  pdftitle={CSB-502: Net Centric Computing -- BHU PYQ}}

\pagestyle{fancy}
\fancyhf{}
\renewcommand{\headrulewidth}{0.4pt}
\renewcommand{\footrulewidth}{0.4pt}
\fancyhead[L]{\small\textit{Banaras Hindu University}}
\fancyhead[R]{\small\textit{CSB-502: Net Centric Computing}}
\fancyfoot[C]{\small Page \thepage\ of \pageref{LastPage}}

\lstset{
  basicstyle=\small\ttfamily,
  frame=single,
  breaklines=true,
  columns=flexible,
  keepspaces=true
}

\newlist{parts}{enumerate}{2}
\setlist[parts,1]{label=(\alph*),leftmargin=*,itemsep=4pt,topsep=3pt}
\setlist[parts,2]{label=(\roman*),leftmargin=*,itemsep=2pt,topsep=2pt}
\newcommand{\pts}[1]{\hfill\textit{[#1]}}

\begin{document}
\begin{center}
  {\large\bfseries BANARAS HINDU UNIVERSITY}\\[2pt]
  {\normalsize B.Sc. (Hons.) Semester V Examination, 2013-14}\\[6pt]
  \rule{0.8\linewidth}{0.5pt}\\[6pt]
  {\large\bfseries Paper: CSB-502 --- Net Centric Computing}\\[4pt]
  {\normalsize Time: Three hours \hfill Maximum Marks: 70}
\end{center}

\rule{\linewidth}{0.4pt}

\smallskip
\noindent\textit{Instructions: ANSWER FIVE QUESTIONS INCLUDING QUESTION NO. 1 WHICH IS COMPULSORY.}

\bigskip

\medskip\noindent\textbf{Question 1.} a) 	extasciitilde{} Whatis mean by flow control? 2
\begin{parts}
  \item = Whatare the fields of UDP header ? 3
  \item Distinguish between bus topology and ring topology. 3
  \item What is count to infinity problem? 3
  \item | Whatare the disadvantages of DES algorithm? 3
\end{parts}

\medskip\noindent\textbf{Question 2.} a) Explain the four basic Network topologies, and cite an advantage of each type. 5
\begin{parts}
  \item What are the responsibilities of the network layer in the Internet model? 5
  \item Describe the multiplexing techniques used to combine digital signals. 4
\end{parts}

\medskip\noindent\textbf{Question 3.} a) Explain in detail, with example, the various error detecting methods. 7
\begin{parts}
  \item Suppose that we model “A Simplex Protocol for a Noisy Channel” using the finite state
  machine model. How many states exist for the communication channel? How many states
  exist for the complete svstem? 7
\end{parts}

\medskip\noindent\textbf{Question 4.} a) What is internetworking? Discuss the various global addressing schemes and the issues in
forwarding the IP (IP packets). ,
\begin{parts}
  \item Explain why routing is very important in Networking. Llustrate any one of the routing
  algorithm used in Network. 7
\end{parts}

\medskip\noindent\textbf{Question 5.} a) Differentiate connectionless and connection oriented services at Transport Layer. Draw the
header format of TCP and explain the purpose of each field. 7
\begin{parts}
  \item What is the difference between open loop congestion control and closed loop congestion
  control? Name the policies that can prevent congestion. 7
\end{parts}

\medskip\noindent\textbf{Question 6.} a) | Whatis the difference between a primary server and a secondary server? What are the two
main categories of DNS messages? 4
\begin{parts}
  \item Write short notes on: 6
  \item WWW ii) browser ili) URL
  \item | What is the mechanism of cookies? 4
\end{parts}

\medskip\noindent\textbf{Question 7.} a) Describe Simple Network Management Protocol (SNMP) and explain the role of MIB
(Management of Information Base). 7
\begin{parts}
  \item Inasymmetric key cryptography, how do you think two persons can establish two pairs of
  keys between themselves? Explain with an example. 7
  RRR
\end{parts}

\vfill
\rule{\linewidth}{0.4pt}
\begin{center}\small
\href{https://bhu-exam.vercel.app/}{Click here to visit our website}\\[4pt]\href{https://me-aryan.vercel.app/}{Click here to visit our contributor}\\[4pt]\href{https://quashan-me.vercel.app/}{Click here to visit our contributor}
\end{center}
\end{document}